\section{AMC 12 12B 2013}
        \begin{enumerate}
        	\item (\textit{AMC 12 12B 2013}) On a particular January day, the high temperature in Lincoln, Nebraska, was $16$ degrees higher than the low temperature, and the average of the high and low temperatures was $3$. In degrees, what was the low temperature in Lincoln that day?


 $\textbf{(A)}\ -13 \qquad \textbf{(B)}\ -8 \qquad \textbf{(C)}\ -5 \qquad \textbf{(D)}\ -3 \qquad \textbf{(E)}\ 11$

 \textbf{Answer: }C



\textbf{Solution}

Let $L$ be the low temperature. The high temperature is $L+16$. The average is $\frac{L+(L+16)}{2}=3$. Solving for $L$, we get $L=\boxed{\textbf{(C)} \  -5}$



	\item (\textit{AMC 12 12B 2013}) Mr. Green measures his rectangular garden by walking two of the sides and finds that it is $15$ steps by $20$ steps. Each of Mr. Green’s steps is $2$ feet long. Mr. Green expects a half a pound of potatoes per square foot from his garden. How many pounds of potatoes does Mr. Green expect from his garden?


 $\textbf{(A)}\ 600 \qquad \textbf{(B)}\ 800 \qquad \textbf{(C)}\ 1000 \qquad \textbf{(D)}\ 1200 \qquad \textbf{(E)}\ 1400$

 \textbf{Answer: }A



\textbf{Solution}

Since each step is $2$ feet, his garden is $30$ by $40$ feet.  Thus, the area of $30(40) = 1200$ square feet.  Since he is expecting $\frac{1}{2}$ of a pound per square foot, the total amount of potatoes expected is $1200 \times \frac{1}{2} =  \boxed{\textbf{(A) }600}$



	\item (\textit{AMC 12 12B 2013}) When counting from $3$ to $201$, $53$ is the $51^{\text{st}}$ number counted. When counting backwards from $201$ to $3$, $53$ is the $n^{\text{th}}$ number counted. What is $n$?


 $\textbf{(A)}\ 146 \qquad \textbf{(B)}\ 147 \qquad \textbf{(C)}\ 148 \qquad \textbf{(D)}\ 149 \qquad \textbf{(E)}\ 150$

 \textbf{Answer: }D



\textbf{Solution}

Note that $n$ is equal to the number of integers between $53$ and $201$, inclusive. Thus, $n=201-53+1=\boxed{\textbf{(D)}\ 149}$



	\item (\textit{AMC 12 12B 2013}) Ray's car averages $40$ miles per gallon of gasoline, and Tom's car averages $10$ miles per gallon of gasoline. Ray and Tom each drive the same number of miles. What is the cars' combined rate of miles per gallon of gasoline?



 $\textbf{(A)}\ 10 \qquad \textbf{(B)}\ 16 \qquad \textbf{(C)}\ 25 \qquad \textbf{(D)}\ 30 \qquad \textbf{(E)}\ 40$

 \textbf{Answer: }B



\textbf{Solution 1}

Let Ray and Tom drive 40 miles. Ray's car would require $\frac{40}{40}=1$ gallon of gas and Tom's car would require $\frac{40}{10}=4$ gallons of gas. They would have driven a total of $40+40=80$ miles, on $1+4=5$ gallons of gas, for a combined rate of $\frac{80}{5}=$ $\boxed{\textbf{(B) }16}$




\textbf{Solution 2}

Taking the harmonic mean of the two rates, we get 
 \[\left(\frac{40^{-1} + 10^{-1}}{2}\right)^{-1} = \frac{2}{\frac{1}{40}+\frac{1}{10}} = \frac{2}{\frac{5}{40}} = \frac{2}{\frac{1}{8}} = \boxed{\textbf{(B) }16}.\]

 -Solution by Joeya

 



\textbf{Solution 3}

Let the number of miles that Ray and Tom each drive be denoted by $m$. Thus the number of gallons Ray's car uses can be represented by $\frac{m}{40}$ and the number of gallons that Tom's car uses can likewise be expressed as $\frac{m}{10}$. Thus the total amount of gallons used by both cars can be expressed as $\frac{m}{40} + \frac{m}{10} = \frac{m}{8}$. The total distance that both drive is equal to $2m$ so the total miles per gallon can be expressed as $\frac{2m}{\frac{m}{8}} = \boxed{\textbf{(B) }16}$



	\item (\textit{AMC 12 12B 2013}) The average age of $33$ fifth-graders is $11$. The average age of $55$ of their parents is $33$. What is the average age of all of these parents and fifth-graders?


 $\textbf{(A)}\ 22 \qquad \textbf{(B)}\ 23.25 \qquad \textbf{(C)}\ 24.75 \qquad \textbf{(D)}\ 26.25 \qquad \textbf{(E)}\ 28$

 \textbf{Answer: }C



\textbf{Solution}

The sum of the ages of the fifth graders is $33 * 11$, while the sum of the ages of the parents is $55 * 33$. Therefore, the total sum of all their ages must be $2178$, and given $33 + 55 = 88$ people in total, their average age is $\frac{2178}{88} = \frac{99}{4} = \boxed{\textbf{(C)}\ 24.75}$.

 


	\item (\textit{AMC 12 12B 2013}) Real numbers $x$ and $y$ satisfy the equation $x^2 + y^2 = 10x - 6y - 34$. What is $x + y$?


 $\textbf{(A)}\ 1 \qquad \textbf{(B)}\ 2 \qquad \textbf{(C)}\ 3 \qquad \textbf{(D)}\ 6 \qquad \textbf{(E)}\ 8$

 \textbf{Answer: }B



\textbf{Solution 1}

If we move every term dependent on $x$ or $y$ to the LHS, we get $x^2 - 10x + y^2 + 6y = -34$. Adding $34$ to both sides, we have $x^2 - 10x + y^2 + 6y + 34 = 0$. We can split the $34$ into $25$ and $9$ to get $(x - 5)^2 + (y + 3)^2 = 0$. Notice this is a circle with radius $0$, which only contains one point. So, the only point is $(5, -3)$, so the sum is $5 + (-3) = 2 \implies \boxed{\textbf{(B)}}$. ~ asdf334

 



\textbf{Solution 2}

If we move every term including $x$ or $y$ to the LHS, we get 
 \[x^2 - 10x + y^2 + 6y = -34.\]  We can complete the square to find that this equation becomes 
 \[(x - 5)^2 + (y + 3)^2 = 0.\]  Since the square of any real number is nonnegative, we know that the sum is greater than or equal to $0$.  Equality holds when the value inside the parhentheses is equal to $0$.  We find that 
 \[(x,y) = (5,-3)\] and the sum we are looking for is 
 \[5+(-3)=2 \implies \boxed{\textbf{(B)}}.\]  - Honestly

 


	\item (\textit{AMC 12 12B 2013}) Jo and Blair take turns counting from $1$ to one more than the last number said by the other person. Jo starts by saying $``1"$, so Blair follows by saying $``1, 2"$. Jo then says $``1, 2, 3"$, and so on. What is the $53^{\text{rd}}$ number said?


 $\textbf{(A)}\ 2 \qquad \textbf{(B)}\ 3 \qquad \textbf{(C)}\ 5 \qquad \textbf{(D)}\ 6 \qquad \textbf{(E)}\ 8$

 \textbf{Answer: }E



\textbf{Solution}

We notice that the number of numbers said is incremented by one each time; that is, Jo says one number, then Blair says two numbers, then Jo says three numbers, etc.  Thus, after nine "turns", $1+2+3+4+5+6+7+8+9=45$ numbers have been said.  In the tenth turn, the eighth number will be the 53rd number said, because $53-45=8$.  Since we are starting from 1 every turn, the 53rd number said will be $\boxed{\textbf{(E) }8}$.

 \textbf{Faster Solution} We notice that the number of numbers said is incremented by one each time; that is, Jo says one number, then Blair says two numbers, then Jo says three numbers, etc.  We notice that the number of numbers is $1 + 2 + 3 + 4 ...$ every time we finish a "turn" we notice the sum of these would be the largest number $\frac{n(n+1)}{2}$ under 53, we can easily see that if we double this it's $n^2 + n \simeq 106$, and we immediately note that 10 is too high, but 9 is perfect, meaning that at 9, 45 numbers have been said so far, $\frac{9(9+1)}{2} = 45$ and $53 - 45 = \boxed{\textbf{(E) }8}$




\textbf{Solution 3}

Let $T(n)$ denote the $n$th triangle number. Then, observe that the $T(n)$th number said is $n$. It follows that the $55$th number is $10$ (as $55 = T(10)$). Thus, the $53$rd number is $10 - 2 = 8$, which is answer choice $\boxed{\textbf{(E)}}$.

 



\textbf{Solution 4}

Let's see when the number $1$ is said: the $1$st number, the $2$nd number, the $4$th number, the $7$th number, and so on. We can see that the $1+\dfrac{n(n+1)}2$th number said is always $1$. That means we must have the $1+\dfrac{9(9+1)}2=46$th number said be $1$. The next $1$ said would then be the $1+\dfrac{10(10+1)}2=56$th number, and $56>53$, so that means every number strictly between the $46$th and $56$th number is not $1$, so let's just count to the $53$rd number!
 \begin{align*}     46\text{th number}&\rightarrow 1\\     47\text{th number}&\rightarrow 2\\     48\text{th number}&\rightarrow 3\\     49\text{th number}&\rightarrow 4\\     50\text{th number}&\rightarrow 5\\     51\text{st number}&\rightarrow 6\\     52\text{nd number}&\rightarrow 7\\     53\text{rd number}&\rightarrow 8, \end{align*}
so the $53$rd number said is $\boxed{\textbf{(E)}\ 8}$.

 


	\item (\textit{AMC 12 12B 2013}) Line $l_1$ has equation $3x - 2y = 1$ and goes through $A = (-1, -2)$. Line $l_2$ has equation $y = 1$ and meets line $l_1$ at point $B$. Line $l_3$ has positive slope, goes through point $A$, and meets $l_2$ at point $C$. The area of $\triangle ABC$ is $3$. What is the slope of $l_3$?


 $\textbf{(A)}\ \frac{2}{3} \qquad \textbf{(B)}\ \frac{3}{4} \qquad \textbf{(C)}\ 1 \qquad \textbf{(D)}\ \frac{4}{3} \qquad \textbf{(E)}\ \frac{3}{2}$

 \textbf{Answer: }B



\textbf{Solution 1}

Line $l_1$ has the equation $y=3x/2-1/2$ when rearranged. Substituting $1$ for $y$, we find that line $l_2$ will meet this line at point $(1,1)$, which is point $B$. We call $\overline{BC}$ the base and the altitude from $A$ to the line connecting $B$ and $C$, $y=-1$, the height. The altitude has length $|-2-1|=3$, and the area of $\triangle{ABC}=3$. Since $A={bh}/2$, $b=2$. Because $l_3$ has positive slope, it will meet $l_2$ to the right of $B$, and the point $2$ to the right of $B$ is $(3,1)$. $l_3$ passes through $(-1,-2)$ and $(3,1)$, and thus has slope $\frac{|1-(-2)|}{|3-(-1)|}=$ $\boxed{\textbf{(B) }\frac{3}{4}}$.

 



\textbf{Solution 2 - Shoelace Theorem}

We know lines $l_1$ and $l_2$ intersect at $B$, so we can solve for that point:
 \[3x-2y=1\]Because $y = 1$ we have: 
 \[3x-2(1) = 1\]
 \[3x-2=1\]
 \[3x=3\]
 \[x = 1\]

 Thus we have $B = (1,1)$. 

 We know that the area of the triangle is $3$, so by Shoelace Theorem we have: 

 
 \[A = \dfrac{1}{2} |(-2x+y-1) - (-2+x-y)|\]
 \[A = \dfrac{1}{2} |-2x+y-1+2-x+y|\]
 \[3 = \dfrac{1}{2} |-2x+y-1+2-x+y|\]
 \[6 = |-3x+2y+1|.\]

 Thus we have two options: 

 
 \[6 = -3x+2y+1\]
 \[5 = -3x+2y\]

 or

 
 \[6 = 3x-2y-1\]
 \[7 = 3x-2y.\]

 Now we must just find a point that satisfies $m_{l_3}$ is positive. 

 Since ${l_3}$ intersects ${l_2}$ at $y=1$, from the second equation: 

 
 \[7 = 3x-2(1)\]
 \[7+2 = 3x\]
 \[x=3\]

 so a valid point here is $(3,1)$. When calculated, the slope of $l_3$ in this situation yields $\boxed{\textbf{(B) }\frac{3}{4}}$.

 


	\item (\textit{AMC 12 12B 2013}) What is the sum of the exponents of the prime factors of the square root of the largest perfect square that divides $12!$ ?


 $\textbf{(A)}\ 5 \qquad \textbf{(B)}\ 7 \qquad \textbf{(C)}\ 8 \qquad \textbf{(D)}\ 10 \qquad \textbf{(E)}\ 12$

 \textbf{Answer: }C



\textbf{Solution}

Looking at the prime numbers under $12$, we see that there are $\left\lfloor\frac{12}{2}\right\rfloor+\left\lfloor\frac{12}{2^2}\right\rfloor+\left\lfloor\frac{12}{2^3}\right\rfloor=6+3+1=10$ factors of $2$, $\left\lfloor\frac{12}{3}\right\rfloor+\left\lfloor\frac{12}{3^2}\right\rfloor=4+1=5$ factors of $3$, and $\left\lfloor\frac{12}{5}\right\rfloor=2$ factors of $5$. All greater primes are represented once or none in $12!$, so they cannot be part of the square. Since we are looking for a perfect square, the exponents on its prime factors must be even, so we can only use $4$ of the $5$ factors of $3$. Therefore, the prime factorization of the square is $2^{10}\cdot3^4\cdot5^2$. To find the square root of this, we halve the exponents, leaving $2^5\cdot3^2\cdot5$. The sum of the exponents is $\boxed{\textbf{(C) }8}$



	\item (\textit{AMC 12 12B 2013}) Alex has $75$ red tokens and $75$ blue tokens. There is a booth where Alex can give two red tokens and receive in return a silver token and a blue token, and another booth where Alex can give three blue tokens and receive in return a silver token and a red token. Alex continues to exchange tokens until no more exchanges are possible. How many silver tokens will Alex have at the end?


 $\textbf{(A)}\ 62 \qquad \textbf{(B)}\ 82 \qquad \textbf{(C)}\ 83 \qquad \textbf{(D)}\ 102 \qquad \textbf{(E)}\ 103$

 \textbf{Answer: }E



\textbf{Solution 1}

If Alex goes to the red booth 3 times, then goes to the blue booth once, Alex can exchange 6 red tokens for 4 silver tokens and one red token. Similarly, if Alex goes to the blue booth 2 times, then goes to the red booth once, Alex can exchange 6 blue tokens for 3 silver tokens and one blue token. Let's call the first combination Combo 1, and the second combination Combo 2.

 In other words, Alex can exchange 5 red tokens for 4 silver tokens as long as he has at least 6 red tokens, and Alex can exchange 5 blue tokens for 3 silver tokens as long as he has at least 6 blue tokens. Hence after performing 14 Combo 1's and 14 Combo 2's, we end up with 5 red, 5 blue, and 98 silver tokens.

 Finally, Alex can visit the blue booth once, then do Combo 1 to end up with 1 red token, 2 blue tokens, and $\boxed{\textbf{(E)}\ 103}$ silver tokens, at which point it is clear he cannot use the booths anymore.

 



\textbf{Solution 2}

We can approach this problem by assuming he goes to the red booth first. You start with $75 \text{R}$ and $75 \text{B}$ and at the end of the first booth, you will have $1 \text{R}$ and $112 \text{B}$ and $37 \text{S}$. We now move to the blue booth, and working through each booth until we have none left, we will end up with:$1 \text{R}$, $2 \text{B}$ and $103 \text{S}$. So, the answer is $\boxed{\textbf{(E)}103}$




\textbf{Solution 3}

Let $x$ denote the number of visits to the first booth and $y$ denote the number of visits to the second booth. Then we can describe the quantities of his red and blue coins as follows:
 \[R(x,y)=-2x+y+75\]
 \[B(x,y)=x-3y+75\]There are no legal exchanges when he has fewer than $2$ red coins and fewer than $3$ blue coins, namely when he has a red coin and $2$ blue coins. We can then create a system of equations:
 \[1=-2x+y+75\]
 \[2=x-3y+75\]Solving yields $x=59$ and $y=44$. Since he gains one silver coin per visit to each booth, he has $x+y=44+59=\boxed{\textbf{(E)}103}$ silver coins in total.

 



\textbf{Solution 4 (Invariance)}

Suppose we set the values of the coins: Red = 4, Blue = 3, Silver = 5. This way, whenever an exchange is made, Alex's total value in coins does not change. His total value at the beginning (and thus also at the end) is $3\cdot 75+4\cdot 75=525$, so his theoretical possible total (if everything works out perfectly) is 105.

 Note that the maximum possible leftover red coins is 1, and the maximum possible leftover blue coins is 2, so the maximum leftover non-silver value is 10. This perfectly gives us 103 silver coins.

 


	\item (\textit{AMC 12 12B 2013}) Two bees start at the same spot and fly at the same rate in the following directions. Bee $A$ travels $1$ foot north, then $1$ foot east, then $1$ foot upwards, and then continues to repeat this pattern. Bee $B$ travels $1$ foot south, then $1$ foot west, and then continues to repeat this pattern. In what directions are the bees traveling when they are exactly $10$ feet away from each other?


 $\textbf{(A)}\ A$ east, $B$ west
$\textbf{(B)}\ A$ north, $B$ south
$\textbf{(C)}\ A$ north, $B$ west
$\textbf{(D)}\ A$ up, $B$ south
$\textbf{(E)}\ A$ up, $B$ west



 
 \textbf{Answer: }A



\textbf{Solution}

Let A and B begin at $(0,0,0)$. In $6$ steps, $A$ will have done his route twice, ending up at $(2,2,2)$, and $B$ will have done his route three times, ending at $(-3,-3,0)$. Their distance is $\sqrt{(2+3)^2+(2+3)^2+2^2}=\sqrt{54} < 10$. We now move forward one step at a time until they are ten feet away:

 7 steps: $A$ moves north to $(2,3,2)$, $B$ moves south to $(-3,-4,0)$, distance of $\sqrt{(2+3)^2+(3+4)^2+2^2}=\sqrt{78} < 10$

 8 steps: $A$ moves east to $(3,3,2)$, $B$ moves west to $(-4,-4,0)$, distance of $\sqrt{(3+4)^2+(3+4)^2+2^2}=\sqrt{102}>10$

 Thus they reach $10$ feet away when $A$ is moving east and B is moving west, between moves 7 and 8. Thus the answer is  $\boxed{\textbf{(A)}}$.

 


	\item (\textit{AMC 12 12B 2013}) Cities $A$, $B$, $C$, $D$, and $E$ are connected by roads $\widetilde{AB}$, $\widetilde{AD}$, $\widetilde{AE}$, $\widetilde{BC}$, $\widetilde{BD}$, $\widetilde{CD}$, and $\widetilde{DE}$. How many different routes are there from $A$ to $B$ that use each road exactly once? (Such a route will necessarily visit some cities more than once.)

                    \begin{figure}[H]
                        \centering
                        \includegraphics[width=0.5\textwidth]{Problem/AMC_12_12B_2013_12_Problem_Figure_1.png}
                    \end{figure}
                    
$\textbf{(A)}\ 7 \qquad \textbf{(B)}\ 9 \qquad \textbf{(C)}\ 12 \qquad \textbf{(D)}\ 16 \qquad \textbf{(E)}\ 18$

 \textbf{Answer: }D



\textbf{Solution}

Note that cities $C$ and $E$ can be removed when counting paths because if a path goes in to $C$ or $E$, there is only one possible path to take out of cities $C$ or $E$.So the diagram is as follows:

 
                    \begin{figure}[H]
                        \centering
                        \includegraphics[width=0.5\textwidth]{Solution/AMC_12_12B_2013_12_Solution_Figure_1.png}
                    \end{figure}
                    

 Now we proceed with casework. Remember that there are two ways to travel from $A$ to $D$, $D$ to $A$, $B$ to $D$ and $D$ to $B$.:

 Case 1 $A \Rightarrow D$: From $D$, if the path returns to $A$, then the next path must go to $B\Rightarrow D \Rightarrow B$. There are $2 \cdot 1 \cdot 2 = 4$ possibilities of the path $ADABDB$. If the path goes to $D$ from $B$, then the path must continue with either $BDAB$ or $BADB$. There are $2 \cdot 2 \cdot 2 = 8$ possibilities. So, this case gives $4+8=12$ different possibilities.

 Case 2 $A \Rightarrow B$: The path must continue with $BDADB$. There are $2 \cdot 2 = 4$ possibilities for this case. 

 Putting the two cases together gives $12+4 = \boxed{\textbf{(D)} \ 16}$

 \textbf{Alcumus Solution (Directly Copied from Alcumus)} Solution:The presence of cities $C$ and $E$ is irrelevant to the problem, because upon entering either city, there is only one road going out. Therefore, we can remove those cities, and instead note that there are two roads connecting $A$ and $D,$ two roads connecting $B$ and $D,$ and one road connecting $A$ and $B.$ We can assume that the order in which each pair of roads is traversed does not matter, and then multiply by $2 \cdot 2 =4$ at the end.

 Now, take cases on whether $B$ or $D$ is visited first:Suppose $D$ is visited first. If the other road back to $A$ is then taken, then the only possibility is to travel to $B$ and then travel the two roads between $B$ and $D$ in either order. If, instead, one of the roads to $B$ is taken, then either $A, D, B$ must be visited in that order, or $D, A, B$ must be visited in that order. This gives $3$ possible routes in total.Suppose $B$ is visited first. Then $D, A, D, B$ must be visited in that order, so there is only one possible route.Putting the two cases together and multiplying by $4$ gives the answer, $4(1+3) = \boxed{16}.$



	\item (\textit{AMC 12 12B 2013}) The internal angles of quadrilateral $ABCD$ form an arithmetic progression. Triangles $ABD$ and $DCB$ are similar with $\angle DBA = \angle DCB$ and $\angle ADB = \angle CBD$. Moreover, the angles in each of these two triangles also form an arithmetic progression. In degrees, what is the largest possible sum of the two largest angles of $ABCD$?


 $\textbf{(A)}\ 210 \qquad \textbf{(B)}\ 220 \qquad \textbf{(C)}\ 230 \qquad \textbf{(D)}\ 240 \qquad \textbf{(E)}\ 250$

 \textbf{Answer: }D



\textbf{Solution 1}

Since the angles of Quadrilateral $ABCD$ form an arithmetic sequence, we can assign each angle with the value $a$, $a+d$, $a+2d$, and $a+3d$. Also, since these angles form an arithmetic progression, we can reason out that $(a)+(a+3d)=(a+d)+(a+2d)=180$.

 For the sake of simplicity, lets rename the angles of each similar triangle. Let $\angle ADB = \angle CBD = \alpha$, $\angle DBA = \angle DCB = \beta$, $\angle CDB = \angle BAD = \gamma$.

 Now the four angles of $ABCD$ are $\beta$, $\alpha + \beta$, $\gamma$, and $\alpha + \gamma$.

 As for the similar triangles, we can name their angles $y$, $y+b$, and 

 


	\item (\textit{AMC 12 12B 2013}) Two non-decreasing sequences of nonnegative integers have different first terms. Each sequence has the property that each term beginning with the third is the sum of the previous two terms, and the seventh term of each sequence is $N$. What is the smallest possible value of $N$ ?


 $\textbf{(A)}\ 55 \qquad \textbf{(B)}\ 89 \qquad \textbf{(C)}\ 104 \qquad \textbf{(D)}\ 144 \qquad \textbf{(E)}\ 273$

 \textbf{Answer: }C



\textbf{Solution 1}

Let the first two terms of the first sequence be $x_{1}$ and $x_{2}$ and the first two of the second sequence be $y_{1}$ and $y_{2}$.  Computing the seventh term, we see that $5x_{1} + 8x_{2} = 5y_{1} + 8y_{2}$.  Note that this means that $x_{1}$ and $y_{1}$ must have the same value modulo $8$.  To minimize, let one of them be $0$; \href{https://artofproblemsolving.com/wiki/index.php?title=WLOG}{WLOG}, assume that $x_{1} = 0$.  Thus, the smallest possible value of $y_{1}$ is $8$; and since the sequences are non-decreasing we get $y_{2} \ge 8$.  To minimize, let $y_{2} = 8$.  Thus, $5y_{1} + 8y_{2} = 40 + 64 = \boxed{\textbf{(C) }104}$.

 



\textbf{Solution 2}

WLOG, let $a_i$, $b_i$ be the sequences with $a_1<b_1$. Then 
 \[N=5a_1+8a_2=5b_1+8b_2\] or 
 \[5a_1+8a_2=5(a_1+c)+8(a_2-d)\] for some natural numbers $c$, $d$. Thus $5c=8d$. To minimize $c$ and $d$, we have $(c,d)=(8,5)$, or 
 \[5a_1+8a_2=5(a_1+8)+8(a_2-5).\] To minimize $a_1$ and $b_1$, we have $(a_1,b_1)=(0,0+c)=(0,8)$. Using the same method, since $b_2\ge b_1$, we have $b_2\ge8$. 

 Thus the minimum $N=5b_1+8b_2= \boxed{\textbf{(C) }104}$.

 



\textbf{Solution 3}

Do the same thing as any of the previous solutions to find the 7th term must be in the form $5a+8b$ and $5c+8d$ where $a\neq c$. Then test the options and find the smallest number that can be expressed in these 2 ways is $\boxed{\textbf{(C) }104}$.

 -Andrew2019

 


	\item (\textit{AMC 12 12B 2013}) 
The number $2013$ is expressed in the form 
 

  $2013 = \frac {a_1!a_2!...a_m!}{b_1!b_2!...b_n!}$,

 
where $a_1 \ge a_2 \ge ... \ge a_m$ and $b_1 \ge b_2 \ge ... \ge b_n$ are positive integers and $a_1 + b_1$ is as small as possible. What is $|a_1 - b_1|$?


 $\textbf{(A)}\ 1 \qquad \textbf{(B)}\ 2 \qquad \textbf{(C)}\ 3 \qquad \textbf{(D)}\ 4 \qquad \textbf{(E)}\ 5$

 \textbf{Answer: }B



\textbf{Solution}

The prime factorization of $2013$ is $61\cdot11\cdot3$. To have a factor of $61$ in the numerator and to minimize $a_1,$ $a_1$ must equal $61$. Now we notice that there can be no prime $p$ which is not a factor of $2013$ such that $b_1<p<61,$ because this prime will not be canceled out in the denominator, and will lead to an extra factor in the numerator. The highest prime less than $61$ is $59$, so there must be a factor of $59$ in the denominator. It follows that $b_1 = 59$ (to minimize $b_1$ as well), so the answer is $|61-59| = \boxed{\textbf{(B) }2}$.

 


	\item (\textit{AMC 12 12B 2013}) Let $ABCDE$ be an equiangular convex pentagon of perimeter $1$. The pairwise intersections of the lines that extend the sides of the pentagon determine a five-pointed star polygon. Let $s$ be the perimeter of this star. What is the difference between the maximum and the minimum possible values of $s$?


 $\textbf{(A)}\ 0 \qquad \textbf{(B)}\ \frac{1}{2} \qquad \textbf{(C)}\ \frac{\sqrt{5}-1}{2} \qquad \textbf{(D)}\  \frac{\sqrt{5}+1}{2} \qquad \textbf{(E)}\ \sqrt{5}$

 \textbf{Answer: }A



\textbf{Solution 1}

The five pointed star can be thought of as five triangles sitting on the five sides of the pentagon. Because the pentagon is equiangular, each of its angles has measure $\frac{180^\circ (5-2)}{5}=108^\circ$, and so the base angles of the aforementioned triangles (i.e., the angles adjacent to the pentagon) have measure $180^\circ - 108^\circ = 72^\circ$. The base angles are equal, so the triangles must be isosceles.

 Let one of the sides of the pentagon have length $x_1$ (and the others $x_2, x_3, x_4, x_5$). Then, by trigonometry, the non-base sides of the triangle sitting on that side of the pentagon each has length $\frac{x_1}{2} \sec 72^\circ$, and so the two sides together have length $x_1 \sec 72^\circ$. To find the perimeter of the star, we sum up the lengths of the non-base sides for each of the five triangles to get $(x_1+x_2+x_3+x_4+x_5) \sec 72^\circ = (1) \sec 72^\circ = \sec 72^\circ$ (because the perimeter of the pentagon is $1$). The perimeter of the star is constant, so the difference between the maximum and minimum perimeters is $\boxed{\textbf{(A)} \ 0}$.

 



\textbf{Solution 2}

The extreme case, that results in the minimum and/or maximum, would probably be a pentagon that approaches a degenerate pentagon. However, due to the way the problem is phrased, we know there exists a minimum and maximum; therefore, we can reasonably assume that the star's perimeter is constant, and answer with $\boxed{\textbf{(A)} \ 0}$.

 


	\item (\textit{AMC 12 12B 2013}) Let $a,b,$ and $c$ be real numbers such that 


 
 \[a+b+c=2, \text{ and}\]

 \[a^2+b^2+c^2=12\]


 What is the difference between the maximum and minimum possible values of $c$?


 $\textbf{(A) }2\qquad \textbf{ (B) }\frac{10}{3}\qquad \textbf{ (C) }4 \qquad \textbf{ (D) }\frac{16}{3}\qquad \textbf{ (E) }\frac{20}{3}$

 \textbf{Answer: }D



\textbf{Solution 1}

Note that $a+b= 2-c$. Now, by \href{https://artofproblemsolving.com/wiki/index.php?title=Cauchy-Schwarz}{Cauchy-Schwarz}, we have that $(a^2+b^2) \ge \frac{(2-c)^2}{2}$. Therefore, we have that $\frac{(2-c)^2}{2}+c^2 \le 12$. We then find the roots of $c$ that satisfy equality and find the difference of the roots. This gives the answer, $\boxed{\textbf{(D)} \ \frac{16}{3}}$.

 Note: We arrived at $(a^2+b^2) \ge \frac{(2-c)^2}{2}$ as we consider the two sets of numbers $(a,b)$ and $(1,1)$. From \href{https://artofproblemsolving.com/wiki/index.php?title=Cauchy-Schwarz}{Cauchy-Schwarz}, we have that $(a^2+b^2)(1^2+1^2) \ge (2\cdot 1 -c\cdot 1)^2$. $a+b = 2-c$ which is how we got the second part of the inequality.

 



\textbf{Solution 2}

This is similar to the first solution but is far more intuitive. From the given, we have 
 \[a + b = 2 - c\] 
 \[a^2 + b^2 = 12 - c^2\] This immediately suggests use of the Cauchy-Schwarz inequality. By Cauchy, we have 
 \[2\,(a^2 + b^2) \geq (a + b)^2\] Substitution of the above results and some algebra yields 
 \[3c^2 - 4c - 20 \leq 0\] This quadratic inequality is easily solved, and it is seen that equality holds for $c = -2$ and $c = \frac{10}{3}$. 

 The difference between these two values is  $\boxed{\textbf{(D)} \ \frac{16}{3}}$.

 



\textbf{Solution 3 (no Cauchy-Schwarz)}

From the first equation, we know that $c=2-a-b$. We substitute this into the second equation to find that
 \[a^2+b^2+(2-a-b)^2=12.\]This simplifies to $2a^2+2b^2-4a-4b+2ab=8$, which we can write as the quadratic $a^2+(b-2)a+(b^2-2b-4)=0$. We wish to find real values for $a$ and $b$ that satisfy this equation. Therefore, the discriminant is nonnegative. Hence,
 \[(b-2)^2-4(b^2-2b-4)\ge0,\]or $-3b^2+4b+20\ge 0$. This factors as $-(3b-10)(b+2)\ge 0$. Therefore, $-2\le b\le \frac{10}{3}$, and by symmetry this must be true for $a$ and $c$ as well.

 Now $a=b=2$ and $c=-2$ satisfy both equations, so we see that $c=-2$ must be the minimum possible value of $c$. Also, $c=\frac{10}{3}$ and $a=b=-\frac{2}{3}$ satisfy both equations, so we see that $c=\frac{10}{3}$ is the maximum possible value of $c$. The difference between these is $\frac{10}{3}-(-2)=\frac{16}{3}$, or $\boxed{\textbf{(D)}}$.

 



\textbf{Solution 4 (geometrical approach)}

From the given, we have $a + b = 2 - c$ and $a^2 + b^2 = 12 - c^2$. The first equation is a line with x and y intercepts of $2-c$ and the second equation is a circle centered at the origin with radius $\sqrt{12-c^2}$. Intuitively, if we want to find the minimum / maximum $c$ such that there still exist real solutions, the two graphs of the equations should be tangent.

 Thus, we have that $\sqrt{2} \cdot \sqrt{12-c^2} = 2-c$, which simplifies to $3c^2-4c-20=0$. Solving the quadratic, we get that the values of $c$ for which the two graphs are tangent are $c=-2$ and $c=\frac{10}{3}$. Thus, our answer is $\boxed{\frac{16}{3}}$.

 



\textbf{Solution 5}

Draw the sphere and the plane represented by the two equations in Cartesian space, with the $z$-axis representing $c$. The intersection between the sphere and plane is a circle. We wish to find the point on the circle where $z$ is minimized and the point where $z$ is maximized. Looking at the graph, it is clear by symmetry that $x = y$ when $z$ is maximized or minimized. Thus, we can set $a = b$. This gives us the following system of equations:

 $2a + c = 2$
$2a^2 + c^2 = 12$

 Solving gives $c = \frac{10}{3}, -2$, which are the maximum and minimum values of $c$ respectively. The answer follows from here.

 



\textbf{Solution 6}

We can consider $a$, $b$, and $c$ to be solutions to a cubic equation. Then, given our information, we have $a+b+c=2$ and $ab+ac+bc=-4$, so our cubic equation looks like this: $x^3-2x^2-4x+m$ where $m$ can be any real number. 

 Since this cubic has $3$ real solutions, it must have both a maximum and a minimum (noticed by looking at the graph). Then the greatest solution is maximized when the other two solutions are the same and is minimized when the other two solutions are the same. 

 Thus, equating $a$ and $b$, we have 
 \[2a+c=2\]and
 \[2a^2+c^2=12\]Solving this, we get the quadratic 
 \[3c^2-4c+4=24 \Rightarrow c=\frac{10}{3}, -2\]implying the answer.

 



\textbf{Solution 7}

Subtracting $c$ from the first equation yields $a+b=2-c$. Subtracting $c^2$ from the second equation yields $a^2+b^2=12-c^2$. Thus we have the equations 
 \[\begin{cases}  a+b=2-c\\ a^2+b^2=12-c^2\\ \end{cases}\]Squaring the first equation yields 
 \[a^2+2ab+b^2=4-4c+c^2\] Subtracting the second equation from this one yields 
 \[2ab=2c^2-4c-8\] 
 \[\iff ab=c^2-2c-4\] Thus we have the system of equations 
 \[\begin{cases}  a+b=2-c\\ ab=c^2-2c-4\\ \end{cases}\]We can reverse \href{https://artofproblemsolving.com/wiki/index.php?title=Vieta\%27s\_Formulas}{Vieta's Formulas}, to get that $a$ and $b$ are roots of the quadratic equation(in $x$) 
 \[x^2-(2-c)x+(c^2-2c-4).\] Because we have $a, b\in \mathbb{R}$, we have that the discriminant of this quadratic equation must be nonnegative. The discriminant is 
 \[(2-c)^2-4(c^2-2c-4)=-3c^2+4c+20,\] and it can be factored as $-3(c+2)(c-\frac{10}{2})$. Since we have $-3(c+2)(c-\frac{10}{2})\geq 0$, we must have $(c+2)(c-\frac{10}{2})\leq 0$. If $c<-2$, then we have that $(c+2)(c-\frac{10}{2})>0$. If $c=-2$, then $(c+2)(c-\frac{10}{2})=0$. If $c\in (-2, \frac{10}{3})$, then $(c+2)(c-\frac{10}{2})<0$. If $c=\frac{10}{3}$, then $(c+2)(c-\frac{10}{2})=0$. If $c>\frac{10}{3}$, then $(c+2)(c-\frac{10}{2})>0$. Thus our inequality holds if and only if $c\in \left [-2, \frac{10}{3}\right]$, and the maxmimum value is $\frac{10}{3}$, whilst the minimum value is $-2$. Thus the difference between the maximum and minimum values is $\frac{10}{3}-(-2)=\boxed{\textbf{(D)} \frac{16}{3}}$

 -vsamc

 


	\item (\textit{AMC 12 12B 2013}) Barbara and Jenna play the following game, in which they take turns. A number of coins lie on a table. When it is Barbara’s turn, she must remove $2$ or $4$ coins, unless only one coin remains, in which case she loses her turn. When it is Jenna’s turn, she must remove $1$ or $3$ coins. A coin flip determines who goes first. Whoever removes the last coin wins the game. Assume both players use their best strategy. Who will win when the game starts with $2013$ coins and when the game starts with $2014$ coins?


 $\textbf{(A)}$ Barbara will win with $2013$ coins and Jenna will win with $2014$ coins. 


 $\textbf{(B)}$ Jenna will win with $2013$ coins, and whoever goes first will win with $2014$ coins. 


 $\textbf{(C)}$ Barbara will win with $2013$ coins, and whoever goes second will win with $2014$ coins.


 $\textbf{(D)}$ Jenna will win with $2013$ coins, and Barbara will win with $2014$ coins.


 $\textbf{(E)}$ Whoever goes first will win with $2013$ coins, and whoever goes second will win with $2014$ coins.


 
 \textbf{Answer: }B



\textbf{Solution 1}

We split into 2 cases: 2013 coins, and 2014 coins.

 $\textbf{2013 coins:}$Notice that when there are $5$ coins left, whoever moves first loses, as they must leave an amount of coins the other person can take. If Jenna goes first, she can take $3$ coins. Then, whenever Barbara takes coins, Jenna will take the amount that makes the total coins taken in that round $5$. (For instance, if Barbara takes $4$ coins, Jenna will take $1$). Eventually, since $2010=0 (\text{mod }5)$ it will be Barbara's move with $5$ coins remaining, so she will lose. If Barbara goes first, on each round, Jenna will take the amount of coins that makes the total coins taken on that round $5$. Since $2013=3 (\text{mod }5)$, it will be Barbara's move with $3$ coins remaining, so she will have to take $2$ coins, allowing Jenna to take the last coin. Therefore, Jenna will win with $2013$ coins.

 $\textbf{2014 coins:}$If Jenna moves first, she will take $1$ coin, leaving $2013$ coins, and she wins as shown above. If Barbara moves first, she can take $4$ coins, leaving $2010$. After every move by Jenna, Barbara will then take the number of coins that makes the total taken in that round $5$. Since $2010=0\text{(mod }5)$, it will be Jenna's turn with $5$ coins left, so Barbara will win. In this case, whoever moves first wins.

 Based on this, the answer is $\boxed{\textbf{(B)}}$




\textbf{Solution 2}

We start with some key observations. 

 First, notice that when there are $0 (\text{mod }5)$ coins left, whoever moves first loses, since the other player can keep taking exactly $5$ minus the amount the first player took, eventually removing the last coin.

 Furthermore, notice that if Barbara is going first when there are $1 (\text{mod }5)$ coins left, she will lose, since the Jenna can follow a similar strategy to the $0 (\text{mod }5)$ case. This will leave Barbara with only 1 coin at the end, in which she loses her turn and Jenna wins.

 


 With these in mind, we split into 2 cases: 2013 coins, and 2014 coins.

 
$\textbf{2013 coins:}$

 If Barbara goes first, she can remove either $2$ or $4$ coins. If she removes $2$ coins, then Jenna will remove $1$ coin, so there will be $0 (\text{mod }5)$ left (hence Barbara will lose). Alternatively, if she removes $4$ coins, then Jenna will remove $3$ coins, so there will be $1 (\text{mod }5)$ left (and Barbara will lose again). In both of these cases, Jenna will win.

 
If Jenna goes first, she will remove $3$ coins, so there will be $0 (\text{mod }5)$ coins left. Hence Jenna will win if she goes first.

 
Putting these together, Jenna will always win with 2013 coins.

 


 $\textbf{2014 coins:}$

 If Barbara goes first, she will remove $4$ coins, so there will be $0 (\text{mod }5)$ coins left. Hence Barbara will win.

 
Or, if Jenna goes first, she will remove $3$ coins, so there will be $1 (\text{mod }5)$ coins left. Hence Jenna will win.

 
Based on these, whoever goes first will win with 2014 coins.

 


 Therefore, the answer is $\boxed{\textbf{(B)}}$.

 


	\item (\textit{AMC 12 12B 2013}) In triangle $ABC$, $AB=13$, $BC=14$, and $CA=15$. Distinct points $D$, $E$, and $F$ lie on segments $\overline{BC}$, $\overline{CA}$, and $\overline{DE}$, respectively, such that $\overline{AD}\perp\overline{BC}$, $\overline{DE}\perp\overline{AC}$, and $\overline{AF}\perp\overline{BF}$. The length of segment $\overline{DF}$ can be written as $\frac{m}{n}$, where $m$ and $n$ are relatively prime positive integers. What is $m+n$?


 $\textbf{(A)}\ 18\qquad\textbf{(B)}\ 21\qquad\textbf{(C)}\ 24\qquad\textbf{(D)}\ 27\qquad\textbf{(E)}\ 30$

 \textbf{Answer: }B



\textbf{Solution 1}

Since $\angle{AFB}=\angle{ADB}=90^{\circ}$, quadrilateral $ABDF$ is cyclic. It follows that $\angle{ADE}=\angle{ABF}$, so $\triangle ABF \sim \triangle ADE$ are similar. In addition, $\triangle ADE \sim \triangle ACD$. We can easily find $AD=12$, $BD = 5$, and $DC=9$ using Pythagorean triples. 

 So, the ratio of the longer leg to the hypotenuse of all three similar triangles is $\frac{12}{15} = \frac{4}{5}$, and the ratio of the shorter leg to the hypotenuse is $\frac{9}{15} = \frac{3}{5}$. It follows that $AF=(\frac{4}{5}\cdot 13), BF=(\frac{3}{5}\cdot 13)$. 

 Let $x=DF$. By \href{https://artofproblemsolving.com/wiki/index.php?title=Ptolemy\%27s\_Theorem}{Ptolemy's Theorem}, we have 
 \[13x+\left(5\cdot 13\cdot \frac{4}{5}\right)= 12\cdot 13\cdot \frac{3}{5} \qquad \Leftrightarrow \qquad 13x+52=93.6.\] Dividing by $13$ we get $x+4=7.2\implies x=\frac{16}{5}$ so our answer is $\boxed{\textbf{(B) }21}$.

 



\textbf{Solution 2}

From solution 1, we know that $AD = 12$ and $DC = 9$. Since $\triangle ADC \sim \triangle DEC$, we can figure out that $DE = \frac{36}{5}$. We also know what $AC$ is so we can figure what $AE$ is: $AE = 15 - \frac{27}{5} = \frac{48}{5}$ . Quadrilateral $ABDF$ is cyclic, implying that $\angle{B} + \angle{DFA}$ = 180°. Therefore, $\angle{B} = 180 - \angle{DFA} = \angle{EFA}$, and triangles $\triangle AEF \sim \triangle ADB$. Solving the resulting proportion gives $EF = 4$. Therefore, $DF = ED - EF = \frac{16}{5}$. $m + n = 16 + 5 = 21$ and our answer is $\boxed{\textbf{(B) } 21}$. 

 $\triangle ADC \sim \triangle DEC$ because of $AA \sim$ . $\angle{ADC} = \angle{DEC} = 90°$. Lets say $\angle{ADE} = x$. So  $\angle{EDC} = 90 - x$ and $\angle{DEC} = 180 - 90 - (90 - x) = x$ so $\angle{ADE} = \angle{DEC}$




\textbf{Solution 3}

Let $\angle{ADF} = \theta,$ and $DF = x.$ Noting that $AD$ is the altitude from $A$ to $BC,$ $AD=12,$ $BD = 5,$ and $DC = 9.$ Notice that since $DE$ is perpendicular to $AC$, $\triangle{ADC}$ is similar to $\triangle{AEC}.$ Therefore since $\triangle{ADC}$ is a $3$-$4$-$5$ triangle, $\cos{\theta}=\frac{3}{5}.$Applying Law of Cosines on $\triangle{ADF}$ we have,
 \[AF^2 = 144+x^2 - 2\cdot 12 \cdot x \cdot \cos{\theta} = 144+x^2 - \frac{72x}{5}.\] Next notice that $\cos{\angle{BDF}} = \cos{(90 + \theta)} = -\sin{\theta} = \frac{-4}{5}.$ Now we apply Law of Cosines on $\triangle{BDF},$
 \[BF^2 = 25 + x^2 - 2 \cdot 5 \cdot x \cdot \cos{\angle{BDF}} = 25 + x^2 + 8x.\]Finally using Pythagorean Theorem on $\triangle{ABF}$ we have, 
 \[AB^2 = AF^2 + BF^2.\] 
 \[13^2 = 144+x^2 - \frac{72x}{5} + 25 + x^2 + 8x.\] 
 \[0 = 2x^2 - \frac{32x}{5}.\] Solving for the nonnegative value for $x$ we find $x = DF = \frac{16}{5} \implies \boxed{21}.$




\textbf{Solution 4}

If we draw a diagram as given, but then add point $G$ on $\overline{BC}$ such that $\overline{FG}\perp\overline{BC}$ in order to use the Pythagorean theorem, we end up with similar triangles $\triangle{DFG}$ and $\triangle{DCE}$. Thus, $FG=\frac35x$ and $DG=\frac45x$, where $x$ is the length of $\overline{DF}$. Using the Pythagorean theorem, we now get 
 \[BF = \sqrt{\left(\frac45x+ 5\right)^2 + \left(\frac35x\right)^2}\] and $AF$ can be found out noting that $AE$ is just $\frac{48}5$ through base times height (since $12\cdot 9 = 15 \cdot \frac{36}5$, similar triangles gives $AE = \frac{48}5$), and that $EF$ is just $\frac{36}5 - x$.  From there, 
 \[AF = \sqrt{\left(\frac{36}5 - x\right)^2 + \left(\frac{48}5\right)^2}.\] Now, $BF^2 + AF^2 = 169$, and squaring and adding both sides and subtracting a 169 from both sides gives $2x^2 - \frac{32}5x = 0$, so $x = \frac{16}5$. Thus, the answer is $\boxed{\textbf{(B)}}$.

 



\textbf{Solution 5 (Power of a Point)}

First, we find $BD = 5$, $DC = 9$, and $AD = 12$ via the Pythagorean Theorem or by using similar triangles. Next, because $DE$ is an altitude of triangle $ADC$, $DE = \frac{AD\cdot DC}{AC} = \frac{36}{5}$. Using that, we can use the Pythagorean Theorem and similar triangles to find $EC = \frac{27}{5}$ and $AE = \frac{48}{5}$. 

 Points $A$, $B$, $D$, and $F$ all lie on a circle whose diameter is $AB$. Let the point where the circle intersects $AC$ be $G$. Using power of a point, we can write the following equation to solve for $AG$: 
 \[DC\cdot BC = CG\cdot AC\] 
 \[9\cdot 14 = CG\cdot 15\] 
 \[CG = 126/15\] Using that, we can find $AG = \frac{99}{15}$, and using $AG$, we can find that $GE = 3$.

 We can use power of a point again to solve for $DF$: 
 \[FE\cdot DE = GE\cdot AE\] 
 \[(\frac{36}{5} - DF)\cdot \frac{36}{5} = 3 \cdot \frac{48}{5}\] 
 \[\frac{36}{5} - DF = 4\] 
 \[DF = \frac{16}{5} = \frac{m}{n}\] Thus, $m+n = 16+5 = 21$ $\boxed{\textbf{(B)}}$.

 



\textbf{Solution 6 (Coord Bash)}

If we draw the diagram like above, but make $BC$ the base, we can set $B(0,0), C(14,0)$ and $A(5,12)$, as $AD$ can quickly be verified to be $12$ by Pythagorean triples or similar triangles. Construct $X$ on $BC$ such that $EX \perp BC$. This implies $\triangle ADC \sim \triangle EXC$ as $AD \parallel EX$, and $\angle ACD = \angle ECX$. Also construct $FY$ such that $FY \perp BC$.

 
Line $\overline {AC}$ has a slope of $-\frac{4}{3}$ by slope formula. Since $D = (5,0)$ and $DE \perp AC,$, the equation of $DE = \frac{3}{4}x - \frac{15}{4}$. 

 Furthermore, $F$ can now be expressed as $(x,\frac{3}{4}x - \frac{15}{4}).$ Since we know $AF \perp BF,$ we can solve for $x$ with the perpendicular slope formula like so: 
 \[\frac{\frac{3}{4}x - \frac{63}{4}}{x - 5} \cdot \frac{\frac{3}{4}x - \frac {15}{4}}{x} = -1\] 
 \[\frac{x-21}{x-5} \cdot \frac{x-5}{x} = -\frac{16}{9}\] 
 \[\frac{x-21}{x} = -\frac{16}{9}\] 
 \[9x = 189 -16x\] 
 \[x = \frac{189}{25}\]

 Plugging $\frac{189}{25}$ into $F$, we get $F(\frac{189}{25},\frac{48}{25}).$ Since $FY \perp DY$, we get that $\triangle FYD$ has side lengths of $FY = \frac{48}{25}$ 

 and $DY = BY - BD = BY - 5 = \frac{64}{25}$. 

 Clearly, $\triangle FYD$ is a $3 - 4 - 5$ pythagorean triple, so $DF = \frac{80}{25} = \frac{16}{5}$.

 $16 + 5 = m + n = \boxed{\textbf{(B) }21}$.

 


	\item (\textit{AMC 12 12B 2013}) For $135^\circ < x < 180^\circ$, points $P=(\cos x, \cos^2 x), Q=(\cot x, \cot^2 x), R=(\sin x, \sin^2 x)$ and $S =(\tan x, \tan^2 x)$ are the vertices of a trapezoid. What is $\sin(2x)$?


 $\textbf{(A)}\ 2-2\sqrt{2}\qquad\textbf{(B)} 3\sqrt{3}-6\qquad\textbf{(C)}\ 3\sqrt{2}-5\qquad\textbf{(D)}\ -\frac{3}{4}\qquad\textbf{(E)}\ 1-\sqrt{3}$

 \textbf{Answer: }A



\textbf{Solution}

Let $f,g,h,$ and $j$ be $\cos{x}, \cot{x}, \sin{x},$ and $\tan{x}$, respectively. Then, we have four points $(f,f^2),(g,g^2),(h,h^2),(j,j^2)$, and a pair of lines each connecting two points must be parallel (as we are dealing with a trapezoid). WLOG, take the line connecting the first two points and the line connecting the last two points to be parallel, so that $\frac{g^2-f^2}{g-f} = \frac{j^2-h^2}{j-h}$, or $g+f = j+h$.

 Now, we must find how to match up $\sin, \cos, \tan, \cot$ to $f,g,h,j$ so that the above equation has a solution. On the interval $135^\circ < x < 180^\circ$, we have $\cot x <-1<\cos x<0<\sin x$, and $\cot x <-1<\tan x<0<\sin x$ so the sum of the largest and the smallest is equal to the sum of the other two, namely, $\sin x+\cot x = \cos x+\tan x$.

 Now, we perform some algebraic manipulation to find $\sin (2x)$:

 $\sin x+\cot x = \cos x+\tan x \\ \sin x - \cos x = \tan x - \cot x = (\sin x - \cos x) (\sin x + \cos x) / (\sin x \cos x) \\ \sin x \cos x = \sin x + \cos x \\ (\sin x \cos x)^2 = (\sin x + \cos x)^2 \\ (\sin x \cos x)^2 = \sin^2 x + 2\sin x \cos x + \cos^2 x \\ (\sin x \cos x)^2 - 2\sin x \cos x - 1 =0$

 Solve the quadratic to find $\sin x \cos x = \frac{2 - 2\sqrt{2}}{2}$, so that $\sin(2x) = 2 \sin x \cos x = \boxed{\textbf{(A)} \ 2 - 2\sqrt{2}}$.

 


	\item (\textit{AMC 12 12B 2013}) Consider the set of $30$ parabolas defined as follows: all parabolas have as focus the point $(0,0)$ and the directrix lines have the form $y=ax+b$ with $a$ and $b$ integers such that $a\in \{-2,-1,0,1,2\}$ and $b\in \{-3,-2,-1,1,2,3\}$. No three of these parabolas have a common point. How many points in the plane are on two of these parabolas?


 $\textbf{(A)}\ 720\qquad\textbf{(B)}\ 760\qquad\textbf{(C)}\ 810\qquad\textbf{(D)}\ 840\qquad\textbf{(E)}\ 870$

 \textbf{Answer: }C



\textbf{Solution 1}

Being on two parabolas means having the same distance from the common focus and both directrices. In particular, you have to be on an angle bisector of the directrices, and clearly on the same "side" of the directrices as the focus. So it's easy to see there are at most two solutions per pair of parabolae. Convexity and continuity imply exactly two solutions unless the directrices are parallel and on the same side of the focus.

 So out of $2\dbinom{30}{2}$ possible intersection points, only $2*5*2*\dbinom{3}{2}$ fail to exist. This leaves $870-60=810=\boxed{\textbf{(C)} \ 810}$ solutions.

 



\textbf{Solution 2}

Through similar reasoning as above in Solution I, determine that two parabolas that have a common focus intersect zero times if their directrixes are parallel and the focus lies on the same side of both directrixes, and intersect twice otherwise. Thereby, as each parabola will intersect $30-3 = 27$ other parabolas twice, we see that the answer is 
 \[2 \times \frac{30 \times 27}{2} = \boxed{810}.\]



	\item (\textit{AMC 12 12B 2013}) Let $m>1$ and $n>1$ be integers. Suppose that the product of the solutions for $x$ of the equation

 \[8(\log_n x)(\log_m x)-7\log_n x-6 \log_m x-2013 = 0\]
is the smallest possible integer. What is $m+n$?


 $\textbf{(A)}\ 12\qquad\textbf{(B)}\ 20\qquad\textbf{(C)}\ 24\qquad\textbf{(D)}\ 48\qquad\textbf{(E)}\ 272$

 \textbf{Answer: }A



\textbf{Solution}

Rearranging logs, the original equation becomes

 
 
 \[\frac{8}{\log n \log m}(\log x)^2 - \left(\frac{7}{\log n}+\frac{6}{\log m}\right)\log x - 2013 = 0\]

 
 

 
	By Vieta's Theorem, the sum of the possible values of $\log x$ is $\frac{\frac{7}{\log n}+\frac{6}{\log m}}{\frac{8}{\log n \log m}} = \frac{7\log m + 6 \log n}{8} = \log \sqrt[8]{m^7n^6}$. But the sum of the possible values of $\log x$ is the logarithm of the product of the possible values of $x$. Thus the product of the possible values of $x$ is equal to $\sqrt[8]{m^7n^6}$.

 
 

 
	It remains to minimize the integer value of $\sqrt[8]{m^7n^6}$. Since $m, n>1$, we can check that $m = 2^2$ and $n = 2^3$ work. Thus the answer is $4+8 = \boxed{\textbf{(A)}\ 12}$.

 
 

 
 

 


	\item (\textit{AMC 12 12B 2013}) Bernardo chooses a three-digit positive integer $N$ and writes both its base-$5$ and base-$6$ representations on a blackboard. Later LeRoy sees the two numbers Bernardo has written. Treating the two numbers as base-$10$ integers, he adds them to obtain an integer $S$. For example, if $N=749$, Bernardo writes the numbers $10444$ and $3245$, and LeRoy obtains the sum $S=13,689$. For how many choices of $N$ are the two rightmost digits of $S$, in order, the same as those of $2N$?


 $\textbf{(A)}\ 5\qquad\textbf{(B)}\ 10\qquad\textbf{(C)}\ 15\qquad\textbf{(D)}\ 20\qquad\textbf{(E)}\ 25$

 \textbf{Answer: }E



\textbf{Solution 1}

First, we can examine the units digits of the number base 5 and base 6 and eliminate some possibilities.

 Say that $N \equiv a \pmod{6}$

 also that $N \equiv b \pmod{5}$

 Substituting these equations into the question and setting the units digits of $2N$ and $S$ equal to each other, it can be seen that $b < 5$ (because otherwise $a$ and $b$ will have different parities), and thus $a=b$.$N \equiv a \pmod{6}$,$N \equiv  a \pmod{5}$,$\implies N=a \pmod{30}$,$0 \le a \le 4$

 Therefore, $N$ can be written as  $30x+a$and $2N$ can be written as $60x+2a$ 

 Just keep in mind that $a$ can be one of five choices: $0, 1, 2, 3,$ or $4$, ;Also, we have already found which digits of $a$ will add up into the units digits of $2N$.

 Now, examine the tens digit, $x$ by using $\mod{25}$ and $\mod{36}$ to find the tens digit (units digits can be disregarded because $a=0,1,2,3,4$ will always work)Then we take $N=30x+a$ $\mod{25}$ and $\mod{36}$ to find the last two digits in the base $5$ and $6$ representation.
 \[N \equiv 30x + a \pmod{36}\]
 \[N \equiv 30x + a \equiv 5x + a \pmod{25}\] Both of those must add up to 
 \[2N\equiv60x + 2a \pmod{100}\]

 ($33 \ge x \ge 4$)

 Now, since $a=0,1,2,3,4$ will always work if $x$ works, then we can treat $x$ as a units digit instead of a tens digit in the respective bases and decrease the mods so that $x$ is now the units digit.
 \[N - a \equiv 5x \pmod{6}\] 
 \[N - a \equiv 6x \equiv x \pmod{5}\] 
 \[2N - 2a \equiv 6x \pmod{10}\]

 Say that $x=5m+n$ (m is between 0-6, n is 0-4 because of constraints on x)Then 

 
 \[N \equiv 5m+n \pmod{5}\] 
 \[N \equiv 25m+5n \pmod{6}\] 
 \[2N\equiv30m + 6n \pmod{10}\]

 and this simplifies to 

 
 \[N \equiv n \pmod{5}\] 
 \[N \equiv m+5n \pmod{6}\]
 \[2N\equiv 6n \pmod{10}\]

 From careful inspection, this is true when

 $n=0, m=6$
$n=1, m=6$
$n=2, m=2$
$n=3, m=2$
$n=4, m=4$

 This gives you $5$ choices for $x$, and $5$ choices for $y$, so the answer is $5* 5 = \boxed{\textbf{(E) }25}$




\textbf{Solution 2}

Notice that there are exactly $1000-100=900=5^2\cdot 6^2$ possible values of $N$. This means, in $100\le N\le 999$, every possible combination of $2$ digits will happen exactly once. We know that $N=900,901,902,903,904$ work because $900\equiv\dots00_5\equiv\dots00_6$.

 We know for sure that the units digit will add perfectly every $30$ added or subtracted, because $\text{lcm }(5, 6)=30$. So we only have to care about cases of $N$ every $30$ subtracted. In each case, $2N$ subtracts $6$/adds $4$, $N_5$ subtracts $1$ and $N_6$ adds $1$ for the $10$'s digit.

 
 \[\textbf{5 }\textcolor{red}{\text{ 0}}\text{ 4 3 2 1 0 }\textcolor{red}{\text{4}}\text{ 3 2 1 0 4 3 2 1 0 4 }\textcolor{red}{\text{3 2}}\text{ 1 0 4 3 2 1 0 4 3 2 }\textcolor{red}{\text{1}}\]

 \[\textbf{6 }\textcolor{red}{\text{ 0}}\text{ 1 2 3 4 5 }\textcolor{red}{\text{0}}\text{ 1 2 3 4 5 0 1 2 3 4 }\textcolor{red}{\text{5 0}}\text{ 1 2 3 4 5 0 1 2 3 4 }\textcolor{red}{\text{5}}\]

 \[\textbf{10}\textcolor{red}{\text{ 0}}\text{ 4 8 2 6 0 }\textcolor{red}{\text{4}}\text{ 8 2 6 0 4 8 2 6 0 4 }\textcolor{red}{\text{8 2}}\text{ 6 0 4 8 2 6 0 4 8 2 }\textcolor{red}{\text{6}}\]

 As we can see, there are $5$ cases, including the original, that work. These are highlighted in $\textcolor{red}{\text{red}}$. So, thus, there are $5$ possibilities for each case, and $5\cdot 5=\boxed{\textbf{(E) }25}$.

 



\textbf{Solution 3}

Notice that $N_5$ ranges from $3$ to $5$ digits and $N_6$ ranges from $3$ to $4$ digits.

 So we can let $N_5 = a_1b_1c_1d_1e_1$ and $N_6 = a_2b_2c_2d_2$. 
 \[S = 10000a_1 + 1000(b_1 + a_2) + 100(c_1 + b_2) + 10(d_1 + c_2) + e_1 + d_2.\] To make it easier to work with mod 100, we will let $N = 625a_1 + 125b_1 + 25c_1 + 5d_1 + e_1$ here. So we have:

 
 \[2N \equiv S \pmod{100} \implies 2(625a_1 + 125b_1 + 25c_1 + 5d_1 + e_1) \equiv 10(d_1 + c_2) + e_1 + d_2 \pmod{100}.\]

 \[50(a_1 + b_1 + c_1) + \cancel{10d_1} + 2e_1 \equiv \cancel{10d_1} + 10c_2 + e_1 + d_2 \pmod{100}.\]

 \[50(a_1 + b_1 + c_1) + e_1 \equiv 10c_2 + d_2 \pmod{100}.\] 

 Now note that $50(a_1 + b_1 + c_1)$ and $10c_2$ will always end with zeros, so the only way for this expression to be true is if $e_1 = d_2.$

 \[50(a_1 + b_1 + c_1) + \cancel{e_1} \equiv 10c_2 + \cancel{d_2} \pmod{100}\]

 \[5(a_1 + b_1 + c_1) \equiv c_2 \pmod{10}.\] 

 So $c_2$ can either be $0$ or $5$.

 Now we can express $N$ in its base-six and base-five forms and set them equal to each other: 

 
 \[625a_1 + 125b_1 + 25c_1 + 5d_1 + \cancel{d_2} = 216a_2 + 36b_2 + 6c_2 + \cancel{d_2}\]

 \[5(125a_1 + 25b_1 + 5c_1 + d_1) = 6c_2 + 36(b_2 + 6a_2).\] 

 Since we established $5 \mid c_2$ (remember, $c_2$ can either be $0$ or $5$), for the right hand side to be a multiple of $5$, $5 \mid b_2 + 6a_2.$ So we have seven possible pairs for $a_2$ and $b_2$: 

 
 \[(a_2, b_2) \rightarrow (5,0), (4,1), (3,2), (2,3), (1,4), (0, 5), (0,0).\] But if we had $(0,0),$ then $N_6$ would only have two digits, and $N$ would not be a three digit number. And if we had $a_2 = 5,$ then $N$ would exceed $1000$ and also not be a three digit number. So we actually have five valid pairs: $(4,1), (3,2), (2,3), (1,4), (0, 5).$ Since a number in one base can be expressed in the other base in exactly one way, and the value of $c_2$ is dependent on $a_1, b_1,$ and $c_1,$ there is one solution $(a_1, b_1, c_1, d_1)$ for every $(a_2, b_2)$ in 

 
 \[N_5 =  625a_1 + 125b_1 + 25c_1 + 5d_1 + d_2 = N_6 =  216a_2 + 36b_2 + 6c_2 + d_2.\] 

 And for every pair we have, there are five choices for $e_1 = d_2$ (any integer from $0 - 4$). So the answer is $5 \cdot 5 = \boxed{25}.$




\textbf{Solution 4}

Observe that the maximum possible value of the sum of the last two digits of the base $5$ number and the base $6$ number is $44+55=99$. Let $N \equiv a \pmod {25}$ and $N \equiv b \pmod {36}$. 

 If $a < \frac{25}{2}$, $2N \equiv 2a \pmod {25}$ and if $a > \frac{25}{2}$, $2N \equiv 2a - 25 \pmod {25}$.

 Using the same logic for $b$, if $b < 18$, $2N \equiv 2b \pmod {36}$, and in the other case $2N \equiv 2b - 36 \pmod {36}$.

 We can do four cases:

 Case 1: $a + b = 2a - 25 + 2b - 36 \implies a + b = 61$.

 For this case, there is trivially only one possible solution, $(a, b) = (25, 36)$, which is equivalent to $(a, b) = (0, 0)$. 

 Case 2: $a + b = 2a - 25 + 2b \implies a + b = 25$.

 Note that in this case, $a \geq 13$ must hold, and $b < 18$ must hold. We find the possible ordered pairs to be: $(13, 12), (14, 11), (15, 10), ..., (24, 1)$ for a total of $12$ ordered pairs.

 Case 3: $a + b = 2a + 2b - 36 \implies a + b = 36$.

 Note that in this case, $b \geq 18$ must hold, and $a < 13$ must hold.We find the possible ordered pairs to be: $(24, 12), (25, 11), (26, 10), ..., (35, 1)$ for a total of $12$ ordered pairs.

 Case 4: $a + b = 2a + 2b$.

 Trivially no solutions except $(a, b) = (0, 0)$, which matches the solution in Case 1, which makes this an overcount.

 By CRT, each solution $(a, b)$ corresponds exactly one positive integer in a set of exactly $\text{lcm} (25, 36) = 900$ consecutive positive integers, and since there are $900$ positive integers between $100$ and $999$, our induction is complete, and our answer is $1 + 12 + 12 = \boxed{\textbf{(E) }25}$.

 


 


	\item (\textit{AMC 12 12B 2013}) Let $ABC$ be a triangle where $M$ is the midpoint of $\overline{AC}$, and $\overline{CN}$ is the angle bisector of $\angle{ACB}$ with $N$ on $\overline{AB}$. Let $X$ be the intersection of the median $\overline{BM}$ and the bisector $\overline{CN}$. In addition $\triangle BXN$ is equilateral with $AC=2$. What is $BX^2$?


 $\textbf{(A)}\  \frac{10-6\sqrt{2}}{7} \qquad \textbf{(B)}\ \frac{2}{9} \qquad \textbf{(C)}\ \frac{5\sqrt{2}-3\sqrt{3}}{8} \qquad \textbf{(D)}\ \frac{\sqrt{2}}{6} \qquad \textbf{(E)}\ \frac{3\sqrt{3}-4}{5}$

 \textbf{Answer: }A



\textbf{Solution 1}

Let $BN=x$ and $NA=y$. From the conditions, let's deduct some convenient conditions that seem sufficient to solve the problem.

 


 \textbf{$M$ is the midpoint of side $AC$.}

 This implies that $[ABX]=[CBX]$. Given that angle $ABX$ is $60$ degrees and angle $BXC$ is $120$ degrees, we can use the area formula to get

 
 \[\frac{1}{2}(x+y)x \frac{\sqrt{3}}{2} = \frac{1}{2} x \cdot CX \frac{\sqrt{3}}{2}\]

 So, $x+y=CX$ .....(1)

 


 \textbf{$CN$ is angle bisector.}

 In the triangle $ABC$, one has $BC/AC=x/y$, therefore $BC=2x/y$.....(2)

 Furthermore, triangle $BCN$ is similar to triangle $MCX$, so $BC/CM=CN/CX$, therefore $BC = (CX+x)/CX = (2x+y)/(x+y)$....(3)

 By (2) and (3) and the subtraction law of ratios, we get 

 
 \[BC=2x/y = (2x+y)/(y+x) = y/x\]

 Therefore $2x^2=y^2$, or $y=\sqrt{2}x$. So $BC = 2x/(\sqrt{2}x) = \sqrt{2}$.

 Finally, using the law of cosine for triangle $BCN$, we get

 
 \[2 = BC^2 = x^2 + (2x+y)^2 - x(2x+y) = 3x^2 + 3xy + y^2 = \left(5+3\sqrt{2}\right)x^2\]

 \[x^2 = \frac{2}{5+3\sqrt{2}} = \boxed{\textbf{(A) }\frac{10-6\sqrt{2}}{7}}.\]




\textbf{Solution 2 (Analytic)}


                    \begin{figure}[H]
                        \centering
                        \includegraphics[width=0.5\textwidth]{Solution/AMC_12_12B_2013_24_Solution_Figure_1.jpg}
                    \end{figure}
                    

 
Let us dilate triangle $ABC$ so that the sides of equilateral triangle $BXN$ are all equal to $2.$ The purpose of this is to ease the calculations we make in the problem. Given this, we aim to find the length of segment $AM$ so that we can un-dilate triangle $ABC$ by dividing each of its sides by $AM$. Doing so will make it so that $AM = 1$, as desired, and doing so will allow us to get the length of $BN$, whose square is our final answer. 

 Let $O$ the foot of the altitude from $B$ to $NX.$ On the coordinate plane, position $O$ at $(0, 0)$, and make $NX$ lie on the x-axis. Since points $N$, $X$, and $C$, are collinear, $C$ must also lie on the x-axis. Additionally, since $NX = 2$, $OB = \sqrt{3}$, meaning that we can position point $B$ at $(0, \sqrt{3})$. Now, notice that line $\overline{AB}$ has the equation $y = \sqrt{3}x + \sqrt{3}$ and that line $\overline{BM}$ has the equation $y = -\sqrt{3}x + \sqrt{3}$ because angles $BNX$ and $BXN$ are both $60^{\circ}$. We can then position $A$ at point $(n, \sqrt{3}(n + 1))$ and $C$ at point $(p, 0)$. Quickly note that, because $CN$ is an angle bisector, $AC$ must pass through the point $(0, -\sqrt{3})$.

 We proceed to construct a system of equations. First observe that the midpoint $M$ of $AC$ must lie on $BM$, with the equation $y = -\sqrt{3}x + \sqrt{3}$. The coordinates of $M$ are $\left(\frac{p + n}{2}, \frac{\sqrt{3}}{2}(n + 1)\right)$, and we can plug in these coordinates into the equation of line $BM$, yielding that 
 \[\frac{\sqrt{3}}{2}(n + 1) = -\sqrt{3}(\frac{p + n}{2}) + \sqrt{3} \implies n + 1 = -p - n + 2 \implies p = -2n + 1.\] For our second equation, notice that line $AC$ has equation $y = \frac{\sqrt{3}}{p}x - \sqrt{3}$. Midpoint $M$ must also lie on this line, and we can substitute coordinates again to get 
 \[\frac{\sqrt{3}}{2}(n + 1) = \frac{\sqrt{3}}{p}(\frac{p + n}{2}) - \sqrt{3} \implies n + 1 = \frac{p + n}{p} - 2 \implies n + 1 = \frac{n}{p} - 1\] 
 \[\implies p = \frac{n}{n + 2}.\]

 Setting both equations equal to each other and multiplying both sides by $(n + 2)$, we have that $-2n^2 - 4n + n + 2 = n \implies -2n^2 - 4n + 2 = 0$, which in turn simplifies into $0 = n^2 + 2n - 1$ when dividing the entire equation by $-2.$ Using the quadratic formula, we have that 
 \[n = \frac{-2 \pm \sqrt{4 + 4}}{2} = -1 - \sqrt{2}.\] Here, we discard the positive root since $A$ must lie to the left of the y-axis. Then, the coordinates of $C$ are $(3 + 2\sqrt{2}, 0)$, and the coordinates of $A$ are $(-1 - \sqrt{2}, -\sqrt{6}).$ Seeing that segment $AM$ has half the length of side $AC$, we have that the length of $AM$ is 
 \[\frac{\sqrt{(3 + 2\sqrt{2} - (-1 - \sqrt{2}))^2 + (\sqrt{6})^2}}{2} = \frac{\sqrt{16 + 24\sqrt{2} + 18 + 6}}{2} = \sqrt{10 + 6\sqrt{2}}.\]

 Now, we divide each side length of $\triangle ABC$ by $AM$, and from this, $BN^2$ will equal $\left(\frac{2}{\sqrt{10 + 6\sqrt{2}}}\right)^2 = \frac{2}{5+3\sqrt{2}} = \boxed{\textbf{(A) }\frac{10-6\sqrt{2}}{7}.}$




\textbf{Solution 3}

By some angle-chasing, we find that $\triangle ANC \sim \triangle BXC$. From here, construct a point $D$ on $AC$ such that $\triangle DXC \sim \triangle ANC$. Now, let $BC = a$, which means that $DM = a - 1$ and $AD = 2 - a$, and let $BN = BX = XN = XD = DN = b$. Note that we want to compute $b^2$. Because $\triangle AND \sim \triangle DXM$, we have:

 
 \[\frac{AN}{2-a} = \frac{b}{a-1} \implies AN = \frac{b(2-a)}{(a-1)}\]

 However, we have more similar triangles. In fact, going back to our original pair of similar triangles - $\triangle ANC$ and $\triangle BXC$ - gives us more similarity ratios:

 
 \[\frac{AN}{AC} = \frac{BX}{BC} \implies \frac{\frac{b(2-a)}{(a-1)}}{2} = \frac{b}{a} \implies a = \sqrt{2}\]

 Since we constructed point $D$ such that $DX$ is parallel to $AB$, we now examine trapezoid $ABXD$. From the variables that we already set up, we know that $AB = b + b\sqrt{2}, BX = XD = b$, and $DA = 2 - \sqrt{2}$. Let $X'$ denote the foot of the perpendicular from $X$ to base $AB$ and define $D'$ similarly. 

 Because $\triangle BXX'$ is a $30, 60, 90$ triangle, $XX' = \frac{b\sqrt{3}}{2}$ and $BX' = \frac{b}{2}$. Thus, $D'A = b\sqrt{2} - \frac{b}{2}$ and $DD' = XX' = \frac{b\sqrt{3}}{2}$. By the Pythagorean Theorem on $\triangle ADD'$,

 
 \[\left (b\sqrt{2} - \frac{b}{2} \right)^2 + \left(\frac{b\sqrt{3}}{2} \right)^2 = \left(2 - \sqrt{2} \right)^2 \implies b^2 = \boxed{\textbf{(A) } \frac{10-6\sqrt{2}}{7}}\].

 


 



\textbf{Solution 4}

Since $\triangle BXN$ is equilateral, let's assume the sides of them are all $a$, and denote the length of $XM$ is $m$. Since $CN$ bisects $\angle BCA$, applying the angle bisector theorem and we can get $BC=\frac{a}{m}$;$AN=2m$. Now applying LOC, we can get $(a+2m)^2+(a+m)^2-2(a+2m)(a+m)\cos\frac{\pi}{3}=1$. We get $a^2+3m^2+3am=1$. Now applying the Stewart theorem in $\triangle BAC$, we can find that ${\frac{a^2}{m^2}+(a+2m)^2=2(1+(a+m)^2)}$, after simplifying, we get ${\frac{a^2}{m^2}-a^2+2m^2=2}$. After observation, the main key for this problem is $a^2$, so we can solve $a$ in term of $m$. Let's see the equation ${\frac{a^2}{m^2}-a^2+2m^2=2}$, we can find that $a=\sqrt{2}m$ so $a^2=2m^2$. Now back solving the first equation we can get that $a=\frac{-3m+\sqrt{4-3m^2}}{m}$cuz the negative one can't work. After solving, we can get that $m^2=\frac{1}{5+3\sqrt{2}}$ so $a^2=2m^2$ and we get $a^2 = \boxed{\textbf{(A) } \frac{10-6\sqrt{2}}{7}}$~ bluesoul

 



\textbf{Solution 5 (Similar Triangles)}

Denote the length of $AN$ as $a$ and the length of $NB$ as $b.$

 Let $M'$ be the midpoint of $\overline{BC} .$ Denote the intersection of $\overline{MM'}$ and $\overline{CN}$ as $X' .$ Note that $MX' = \frac12 AN = \frac{a}{2}$ and $CN = 2\cdot NX' .$ As $\overline{MM'} || \overline{AB},$ we have that $\triangle MXX' \sim \triangle BXN$ or $\triangle MXX'$ is equilateral and $XX' =MX' = \frac{a}{2}.$ Thus, $CN = 2b+a$ and $CX=a+b.$

 Observe that 

 
 \[\triangle BXC\sim \triangle ANC \implies \frac{BX}{XC} =\frac{AN}{NC} \implies \frac{b}{b+a} = \frac{a}{2b+a} \implies a = \sqrt 2 b.\]

 By the angle bisector theorem, we have that $BC=\frac{2b}{a} = \sqrt 2.$

 We apply the Law of Cosines on $\triangle BXC$ as follows:
 \[BC^2=BX^2+XC^2-2BX\cdot XC\cdot \cos120^\circ\]
 \begin{align*} 2&=b^2+(\sqrt 2+1)^2b^2 +(\sqrt2 + 1)b^2 \\ &=b^2(5+3\sqrt 2) \end{align*}
or
 \[\boxed{b^2=\textbf{(A) } \frac{10-6\sqrt{2}}{7}}\]




\textbf{Solution 6}


                    \begin{figure}[H]
                        \centering
                        \includegraphics[width=0.5\textwidth]{Solution/AMC_12_12B_2013_24_Solution_Figure_2.png}
                    \end{figure}
                    $\angle BXC = \angle ANC$, $\angle BCX = \angle ACN$, $\triangle BCX \sim \triangle ACN$, $\frac{CX}{CN} = \frac{BC}{AC}$
$\angle MBC = \angle BAC$, $\angle BCM = \angle ACB$, $\triangle BCM \sim \triangle ACB$, $\frac{BC}{AC} = \frac{CM}{BC}$, $BC = \sqrt{2}$

 Let $BX = x$, $CN = CX + x$, $\frac{CX}{CX + x} = \frac{\sqrt{2}}{2}$, $2CX = CX \sqrt{2} + x \sqrt{2}$, $CX = x(\sqrt{2} + 1)$
$BC^2 = BX^2 + CX^2 - 2 \cdot BX \cdot CX \cdot \cos 120^\circ$
$2 = x^2(\sqrt{2} + 1)^2 + x^2 + x^2(\sqrt{2} + 1) = 5x^2 + 3x^2 \sqrt{2}$
$x^2 = \frac{2}{5 + 3 \sqrt{2} } = \boxed{\textbf{(A) }\frac{10-6\sqrt{2}}{7}}$




\textbf{Solution 7}

Let $BN = m$ and $AN = n$. Let $BC = a$. It's known that $AM = MC = 1$. By the Angle Bisector Theorem:

 $\frac{BN}{AN} = \frac{BC}{AC} \implies \frac{m}{n} = \frac{a}{2} \implies n = \frac{2m}{a}$.

 Now we will apply the Angle Bisector Theorem again on triangle $BCM$. Hence:

 $\frac{BC}{CM} = \frac{BX}{MC}$ Since $BNX$ is equilateral, $BN = BX = m$. Let $MX = k$. Hence we proceed with:

 $\frac{a}{1} = \frac{m}{k} \implies k = \frac{m}{a}$

 We will now use the Law of Sines on triangle $BXC$ by first denoting $\angle BCN = \angle ACN = \theta$. Hence:

 $\frac{\sin(\theta)}{m} = \frac{\sin(\pi - \frac{\pi}{3})}{a}$ since equilateral triangles have radian measures of $\frac{\pi}{3}$ and $N, X, C$ are collinear. We simplify this keeping in mind that $\sin(\pi - \frac{\pi}{3}) = \sin(\frac{2\pi}{3}) = \frac{\sqrt{3}}{2}$. Therefore:

 $\sin(\theta) = \frac{m\sqrt{3}}{2a}$

 We will now find $\cos(2\theta)$. By the double-angle identity:

 $\cos(2\theta) = 1 - 2\sin^{2}(\theta) = 1 - 2(\frac{m\sqrt{3}}{2a})^{2} = 1 - \frac{3m^{2}}{2a^{2}}$.

 We will now use Law of Cosines on triangle $BMC$. We proceed to get:

 $a^{2} + 1 - 2a\cos(2\theta) = (m + \frac{m}{a})^{2}$
$\implies a^{2} + 1 - 2a(1 - \frac{3m^{2}}{2a^{2}}) = m^{2} + \frac{2m^{2}}{a} + \frac{m^{2}}{a^{2}}$
$\implies a^{2} + 1 - 2a + \frac{3m^{2}}{a} = m^{2} + \frac{2m^{2}}{a} + \frac{m^{2}}{a^{2}}$
$\implies m^{2} + \frac{m^{2}}{a^{2}} = a^{2} + 1 - 2a + \frac{m^{2}}{a}$

 
We will now use Law of Cosines again on triangle $ABC$. We proceed to get:

 $a^{2} + 4 - 4a\cos(2\theta) = (m + \frac{2m}{a})^{2}$
$\implies a^{2} + 4 - 4a(1 - \frac{3m^{2}}{2a^{2}}) = m^{2} + \frac{4m^{2}}{a} + \frac{4m^{2}}{a^{2}}$
$\implies a^{2} + 4 - 4a + \frac{6m^{2}}{a} = m^{2} + \frac{4m^{2}}{a} + \frac{4m^{2}}{a^{2}}$
$\implies m^{2} + \frac{4m^{2}}{a^{2}} = a^{2} + 4 - 4a + \frac{2m^{2}}{a}$

 From this new recent equation, we get:

 $m^{2} - a^{2} = 4 - 4a + \frac{2m^{2}}{a} - \frac{4m^{2}}{a^{2}}$

 From the last equation we found, we can get:

 $m^{2} - a^{2} = 1 - 2a + \frac{m^{2}}{a} - \frac{m^{2}}{a^{2}}$

 So we get:

 $4 - 4a + \frac{2m^{2}}{a} - \frac{4m^{2}}{a^{2}} = 1 - 2a + \frac{m^{2}}{a} - \frac{m^{2}}{a^{2}}$
$\implies 3 + \frac{2m^{2}}{a} - \frac{4m^{2}}{a^{2}} = 2a + \frac{m^{2}}{a} - \frac{m^{2}}{a^{2}}$
$\implies 3 = 2a - \frac{m^{2}}{a} + \frac{3m^{2}}{a^{2}}$

 We will use Law of Cosines one last time on triangle $ABM$ noting that $AB = BN + AN = m + \frac{2m}{a}$ and $BM = BX + XM = m + \frac{m}{a}$

 We have:

 $(m + \frac{2m}{a})^{2} + (m + \frac{m}{a})^{2} - 2(m + \frac{2m}{a})(m + \frac{m}{a})\cos(\frac{\pi}{3}) = 1$

 We can simplify this to get:

 $m^{2} + \frac{3m^{2}}{a} + \frac{3m^{2}}{a^{2}} = 1$

 Now we are ready to solve:

 $3 = 2a - \frac{m^{2}}{a} + \frac{3m^{2}}{a^{2}}$ (1)

 $m^{2} + \frac{3m^{2}}{a} + \frac{3m^{2}}{a^{2}} = 1$ (2)

 Note that we can write:

 $\frac{3m^{2}}{a^{2}} = 3 - 2a + \frac{m^{2}}{a}$
$\frac{3m^{2}}{a^{2}} = 1 - m^{2} - \frac{3m^{2}}{a}$

 Hence, $3 - 2a + \frac{m^{2}}{a} = 1 - m^{2} - \frac{3m^{2}}{a} \implies 2a - 2 = \frac{4m^{2}}{a} + m^{2}$ 

 From this new equation we can solve for $m^{2}$ to get:

 $m^{2} = \frac{2a - 2}{\frac{4}{a} + 1}$

 From equation (2), we can also solve for $m^{2}$ to get:

 $m^{2} = \frac{1}{1 + \frac{3}{a} + \frac{3}{a^{2}}}$

 Hence:

 $\frac{2a - 2}{\frac{4}{a} + 1} = \frac{1}{1 + \frac{3}{a} + \frac{3}{a^{2}}}$

 We can clear fractions and cross multiply to get:

 $(2a^{2} - 2a)(a^{2} + 3a + 3) = a^{2}(a + 4) \implies (2a - 2)(a^{2} + 3a + 3) = a(a + 4) = a^{2} + 4a$

 Hence:

 $2a^{3} + 4a^{2} - 6 = a^{2} + 4a \implies 2a^{3} + 3a^{2} - 4a - 6 = 0$

 Magnificently, we can factor this! We get:

 $(a^{2} - 2)(2a + 3) = 0$ following factoring by grouping

 Now, it's clear $a = \pm \sqrt{2}, -\frac{3}{2}$ Since $a$ is a side length, only $a = \sqrt{2}$ makes sense. Remember, the problem asked for $m^{2} = \frac{2a - 2}{\frac{4}{a} + 1} = \frac{2\sqrt{2} - 2}{2\sqrt{2} + 1} = \frac{10 - 6\sqrt{2}}{7}$ so our answer is $(A)$. 

 


	\item (\textit{AMC 12 12B 2013}) Let $G$ be the set of polynomials of the form

 \[P(z)=z^n+c_{n-1}z^{n-1}+\cdots+c_2z^2+c_1z+50,\]
where $c_1,c_2,\cdots, c_{n-1}$ are integers and $P(z)$ has distinct roots of the form $a+ib$ with $a$ and $b$ integers. How many polynomials are in $G$?


 $\textbf{(A)}\ 288\qquad\textbf{(B)}\ 528\qquad\textbf{(C)}\ 576\qquad\textbf{(D)}\ 992\qquad\textbf{(E)}\ 1056$

 \textbf{Answer: }B



\textbf{Solution 1}

If we factor into irreducible polynomials (in $\mathbb{Q}[x]$), each factor $f_i$ has exponent $1$ in the factorization and degree at most $2$ (since the $a+bi$ with $b\ne0$ come in conjugate pairs with product $a^2+b^2$). Clearly we want the product of constant terms of these polynomials to equal $50$; for $d\mid 50$, let $f(d)$ be the number of permitted $f_i$ with constant term $d$. It's easy to compute $f(1)=2$, $f(2)=3$, $f(5)=5$, $f(10)=5$, $f(25)=6$, $f(50)=7$, and obviously $f(d) = 1$ for negative $d\mid 50$.

 Note that by the distinctness condition, the only constant terms $d$ that can be repeated are those with $d^2\mid 50$ and $f(d)>1$, i.e. $+1$ and $+5$. Also, the $+1$s don't affect the product, so we can simply count the number of polynomials with no constant terms of $+1$ and multiply by $2^{f(1)} = 4$ at the end.

 We do casework on the (unique) even constant term $d\in\{\pm2,\pm10,\pm50\}$ in our product. For convenience, let $F(d)$ be the number of ways to get a product of $50/d$ without using $\pm 1$ (so only using $\pm5,\pm25$) and recall $f(-1) = 1$; then our final answer will be $2^{f(1)}\sum_{d\in\{2,10,50\}}(f(-d)+f(d))(F(-d)+F(d))$. It's easy to compute $F(-50)=0$, $F(50)=1$, $F(-10)=f(5)=5$, $F(10)=f(-5)=1$, $F(-2)=f(-25)+f(-5)f(5)=6$, $F(2)=f(25)+\binom{f(5)}{2}=16$, so we get
 \[4 [ (1+3)(6+16) + (1+5)(1+5) + (1+7)(0+1) ] = 4[132] = \boxed{\textbf{(B) }528}\]




\textbf{Solution 2}

Disregard sign; we can tack on $x-1$ if the product ends up being negative.

 $1: \pm i,-1$ (2) (1 is not included)

 $2: \pm 2, \pm 1\pm i$ (4)

 $5: \pm 2\pm i, \pm 1\pm 2i, \pm 5$ (6)

 $10: \pm 3\pm i, \pm 1\pm 3i, \pm 10$ (6)

 $25: \pm 25, \pm 3\pm 4i, \pm 4\pm 3i, \pm 5i$ (7)

 $50: \pm 50, \pm 1\pm 7i, \pm7\pm i, \pm 5\pm 5i$ (8)

 Our answer is $2^2\left(4\cdot\binom{6}{2}+6\cdot 6+4\cdot 7+8\right)=\boxed{528.}$




\textbf{Solution 3}

By Vieta's formula $50$ is the product of all $n$ roots. As the roots are all in the form $a + bi$, there must exist a conjugate $a-bi$ for each root.

 $(a+bi)(a-bi) = a^2 + b^2$
$50 = 2 \cdot 5^2$

 If $a \neq b \neq 0$, the roots can be $a \pm bi$, $-a \pm bi$, $b \pm ai$, $-b \pm ai$, totaling $4$ pairs of roots.

 If $a = b$, the roots can be $a \pm ai$, $-a \pm ai$, totaling $2$ pairs of roots.

 If $a \neq b$, $b = 0$, the roots can be $\pm a$, $\pm ai$, totaling $2$ pairs of roots.

 
 \begin{align*} 2 \cdot 25 &= (1^2+1^2)5^2 &: 2 \cdot 2 = 4\\ 2 \cdot 25 &= 2 \cdot 5^2 &: 2 \cdot 2 = 4\\ 2 \cdot 25 &= (1^2+1^2) \cdot (3^2+4^2) &: 2 \cdot 4 = 8\\ 2 \cdot 25 &= 2 \cdot (3^2+4^2) &: 2 \cdot 4 = 8 \end{align*}


 \begin{align*} 10 \cdot 5 &= (1^2+3^2)(1^2+2^2) &&: 4 \cdot 4 = 16\\ 10 \cdot 5 &= 10 \cdot (1^2+2^2) &&: 2 \cdot 4 = 8\\ 10 \cdot 5 &= (1^2+3^2) \cdot 5 &&: 4 \cdot 2 = 8\\ 10 \cdot 5 &= 10 \cdot 5 &&: 2 \cdot 2 = 4\\ \end{align*}


 \begin{align*} 2 \cdot 5  \cdot 5&= (1^2+1^2)(1^2+2^2)(1^2+2^2) &&: 2 \cdot 4 \cdot 4 = 32\\ 2 \cdot 5  \cdot 5&= 2 \cdot (1^2+2^2)(1^2+2^2) &&: 2 \cdot 4 \cdot 4 = 32\\ 2 \cdot 5  \cdot 5&= 2 \cdot 5 \cdot (1^2+2^2) &&: 2 \cdot 2 \cdot 4 = 16\\ 2 \cdot 5  \cdot 5&= 2 \cdot 5 \cdot 5 &&: 2 \cdot 2 \cdot 2 = 8\\ 2 \cdot 5  \cdot 5&= (1^2+1^2) \cdot 5 \cdot (1^2+2^2) &&: 2 \cdot 2 \cdot 4 = 16\\ 2 \cdot 5  \cdot 5&= (1^2+1^2) \cdot 5 \cdot 5 &&: 2 \cdot 2 \cdot 2 = 8\\ \end{align*}


 \begin{align*} (1^2+7^2) &: 4\\ (5^2+5^2) &: 2\\ 50 &: 2 \end{align*}

$4+4+8+8+16+8+8+4+32+32+16+8+16+8+4+2+2 = 176$

 For each case $1^2$ can be added, yielding 2 more cases $(\pm 1, \pm i)$. $176 \cdot 3 = \boxed{\textbf{(B) }528}$



\end{enumerate}